% Default skeleton for research-hub papers (math-paper-writing skill).
% Every TODO must be replaced; check_paper.py fails while any remain.
% Remove these comment lines from the release copy.
\documentclass[11pt,reqno]{amsart}
\usepackage[T1]{fontenc}
\usepackage{amsmath,amssymb,amsthm,lmodern}
\usepackage{mathtools}
\usepackage[margin=1in]{geometry}
\usepackage{microtype}
% \usepackage{tikz-cd} % only if the paper has commutative diagrams
\usepackage{hyperref}
\hypersetup{hidelinks,
 pdftitle={TODO title},
 pdfauthor={TODO author}}

\theoremstyle{plain}
\newtheorem{theorem}{Theorem}[section]
\newtheorem{proposition}[theorem]{Proposition}
\newtheorem{lemma}[theorem]{Lemma}
\newtheorem{corollary}[theorem]{Corollary}
\newtheorem{conjecture}[theorem]{Conjecture}
\newtheorem{question}[theorem]{Question}
% Lettered headline results (Theorem A, B) for the introduction:
\newtheorem{introtheorem}{Theorem}
\renewcommand{\theintrotheorem}{\Alph{introtheorem}}
\theoremstyle{definition}
\newtheorem{definition}[theorem]{Definition}
\newtheorem{example}[theorem]{Example}
\theoremstyle{remark}
\newtheorem{remark}[theorem]{Remark}

\newcommand{\PP}{\mathbb{P}}
\newcommand{\OO}{\mathcal{O}}
\newcommand{\Db}{\mathrm{D}^{\mathrm{b}}}
\DeclareMathOperator{\Hom}{Hom}
\DeclareMathOperator{\Ext}{Ext}
\DeclareMathOperator{\End}{End}
\DeclareMathOperator{\rk}{rk}
\DeclareMathOperator{\coker}{coker}
\DeclareMathOperator{\Pic}{Pic}
\newcommand{\sExt}{\mathcal{E}\!\mathit{xt}}

\begin{document}

\title{TODO specific title naming the object and the result}

% Author block: copy from the user's latest approved preprint and confirm
% with the user. Never invent an affiliation, address or identifier.
\author{TODO author}
\address{TODO address}
\email{TODO email}
\date{}
\subjclass[2020]{Primary TODO; Secondary TODO}
\keywords{TODO keywords}

\begin{abstract}
TODO 80--150 words: state the main result in words with its hypotheses,
then the method if distinctive, then the secondary results. No citations.
\end{abstract}

\maketitle

\section{Introduction}

TODO one or two paragraphs: the problem and why it matters.

\begin{theorem}\label{thm:main}
TODO exact statement, identical to the statement proved in
Section~\ref{sec:proof}.
\end{theorem}

\subsection{Relation to earlier work}
TODO who proved what, in what generality, with locators; what is new here
and which difficulty earlier methods could not handle.

\subsection{Idea of the proof}
TODO one or two paragraphs naming the mechanism; say which step is new.

\subsection{Limits and questions}
TODO examples showing that the hypotheses are needed; what is not proved.

\subsection{Organization}
TODO a roadmap by section number, flagging routine sections and appendices.

\subsection{Conventions}
TODO base field and characteristic, the meaning of ``variety'', derived
functors, and the ordering conventions used (see the skill's
references/algebraic-geometry.md).

\section{TODO preliminaries}\label{sec:prelim}

\section{Proof of Theorem~\ref{thm:main}}\label{sec:proof}

\begin{proof}[Proof of Theorem~\ref{thm:main}]
TODO assemble the lemmas, citing each by number.
\end{proof}

\appendix
\section{TODO routine verifications}\label{app:routine}

\subsection*{Acknowledgements}
TODO the user's decision, including any AI-assistance statement.

\begin{thebibliography}{TODO}
% Only entries matching a dossier row or a hub source record.
% Cite with locators: \cite[Theorem~3.5]{Key}.
\bibitem[TODO]{TODO}
TODO author, \emph{title}, journal \textbf{volume} (year), pages.
\end{thebibliography}

\end{document}
